\begin{lemma}
\label{lemma-pid-jacobson}
The ring $\mathbf{Z}$ is a Jacobson ring.
More generally, let $R$ be a ring such that
\begin{enumerate}
\item $R$ is a domain,
\item $R$ is Noetherian,
\item any nonzero prime ideal is a maximal ideal, and
\item $R$ has infinitely many maximal ideals.
\end{enumerate}
Then $R$ is a Jacobson ring. In (2), it suffices more generally
that $\Spec(R)$ is a Noetherian topological space.
\end{lemma}

\begin{proof}
Assume (1), (3), (4) and that $\Spec(R)$ is Noetherian. The statement
means that $(0) = \bigcap_{\mathfrak m \subset R} \mathfrak m$.
Since $R$ has infinitely many maximal ideals it suffices to
show that any nonzero $x\in R$ belongs to only finitely many
maximal ideals, or equivalently that $V(x)$ is finite.
By Lemma \ref{lemma-spec-closed}
we see that $V(x)$ is homeomorphic to $\Spec(R/xR)$.
By assumption (3) every prime of $R/xR$ is minimal and hence
corresponds to an irreducible component of $\Spec(R/xR)$
(Lemma \ref{lemma-irreducible}).
If $R$ is Noetherian, then $\Spec(R)$ is Noetherian by
Lemma \ref{lemma-Noetherian-topology}. More generally assume only
this topological condition. Its closed subspace $V(x)$ is Noetherian,
and therefore has finitely many irreducible components by Topology,
Lemma \ref{topology-lemma-Noetherian}. Every prime in $V(x)$ is
nonzero and hence maximal by (3), so each irreducible component is
a singleton. Thus $V(x)$ is finite. The infinitely many maximal
ideals cannot all contain this $x$, proving their intersection is zero.
The nonzero primes are already maximal by (3), so every prime is an
intersection of maximal ideals and Lemma \ref{lemma-jacobson-prime}
proves that $R$ is Jacobson.
\end{proof}